\scanpage{43}\mspage{42}
\uline{characterized} by the following conditions :
% slanted
\margin{This is false, it implies (i) (ii) (iii) (iv), but these conditions are by no means sufficient, see below ! see~48}
\margin{i.e $\overline{\tau'}\in\underline{\underline{\overline{\mathscr{T}}}}$ $\left\{\vphantom{\rule{0pt}{4ex}}\right.$}

(i) $\overline{\tau'}$ is finer than $\widetilde{\tau}$, and induces on $J'$ $\widetilde{\tau}\mid J'$.

(ii) $\overline{\tau'}\mid\overset{\deleted{\unsure{\sim|\mathord{-}}}}{U'}=I'\setminus J'$ coincides with $\widetilde{\tau}\mid\overset{\deleted{\sim}}{U'}$

(iii) $\varphi\mid U'$ induces \deleted{a bijection} \inserted{\unsure{preorder isomorphism}} $U'\longrightarrow U=S\setminus T$
% slanted
\margin{\deleted{That is to say}}

(\unsure{iv|iiv}) The topology of $I$ is a final quotient for the maps $I'\mathbin{\text{---}}I$, $J\mathbin{\text{---}}I$ ($I'$ equipped \deleted{that of $I$, \unsure{jointly}} \inserted{with $\overline{\tau'}$, $J$ \deleted{\unsure{with}} equipped with $\overline{\tau}\mid J$}) (Which excludes in particular the trivial solution for (i) (ii), obtained by taking the sum topology on \deleted{$J'$} $I'=J'\underset{\text{disjoint}}{\cup}\overset{\deleted{\sim}}{U'}$, unless $I$ itself is a topological sum $J\sqcup\underset{\substack{\Vert\\I\setminus J}}{U}$).
% slanted
\margin{It is this condition that I had at first forgotten, on page 31 !}
% slanted
\margin{\deleted{That is to} say that $\tau'$ on $S$ is characterized by the analogous properties (i) (ii) (iii) (iv), the last \unsure{meaning|means} that the topology $\text{\textcircled{$\tau$}}$ on $S$ is final for $\unsure{\tau''|\tau'}$ on $S$ and $\tau$ on $T$ (id $S\unsure{\Rightarrow|\rightrightarrows|\to}S$ $T\unsure{\hookrightarrow|\to}S$)}
% slanted
\margin{in fact, it is also independent of the \unsure{ultra-topological} hypothesis}

Proof

\deleted{a) and b).} \inserted{(a b c)} It is clear that the sets pinned in c) are those which satisfy $1^\circ)$, $2^\circ)$ $3^\circ)$ in b), and that the representation in c) is unique. It remains to see, to prove (a,b,c), that the sets form a topology, \deleted{and} that it \deleted{\unsure{has the properties|has the properties}} \inserted{is} in $\underline{\mathscr{T}}$, and is least fine of the topologies $\in\underline{\mathscr{T}}$.

1) These are the closed sets of a topology.

Stable under arbitrary intersections : clear

Stable under finite unions : ditto.

It is even stable under arbitrary unions if this is so for $(T,\tau'_0)$, $(\deleted{\unsure{I|D}}\unsure{S|I},\tau)$, whence the first assertion of \unsure{d)|a)}.
